\chapter[\emj{📑}~Topological Sort (Ordenamiento Topológico)]{\emj{📑}~Topological Sort (Ordenamiento Topológico)}

\begin{dsaquees}

Ponerte la ropa en el orden correcto 👕. No puedes ponerte los zapatos
antes que los calcetines, ni el saco antes que la camisa.

Cada prenda "depende" de otra. El Ordenamiento Topológico te da una
lista válida de en qué orden vestirte para que nunca rompas una regla:
"primero A, luego B".

Solo funciona si NO hay contradicciones (ciclos). Si la regla dijera
"ponte los zapatos antes que los calcetines Y los calcetines antes que
los zapatos", sería imposible: eso es un ciclo.

\end{dsaquees}

\begin{dsaotraforma}

El Topological Sort es un ordenamiento lineal de los vértices de un
grafo dirigido acíclico (DAG) tal que, para toda arista $u \to v$, el
nodo $u$ aparece antes que $v$ en el orden.

Solo existe en grafos dirigidos y SIN ciclos (DAG). Si hay un ciclo,
no hay orden válido. Se usa para modelar dependencias: compilación de
módulos, prerrequisitos de materias, planificación de tareas y
resolución de fórmulas en hojas de cálculo.

Hay dos formas de calcularlo: el algoritmo de Kahn (BFS con in-degree)
y el enfoque DFS (post-orden invertido).

\end{dsaotraforma}

\begin{dsadiagrama}
\begin{verbatim}
Grafo de dependencias (DAG):

  [Calcetines] --> [Zapatos]
  [Camisa] --> [Saco] --> [Corbata]
  [Ropa interior] --> [Pantalon] --> [Zapatos]

In-degree (aristas entrantes):
  Calcetines:0  Camisa:0  RopaInterior:0
  Pantalon:1  Zapatos:2  Saco:1  Corbata:1

Un orden topologico valido:
  RopaInterior -> Pantalon -> Calcetines -> Zapatos
  -> Camisa -> Saco -> Corbata

(No es unico: cualquier orden que respete las flechas sirve).
\end{verbatim}
\end{dsadiagrama}

\begin{dsacomofunciona}
\begin{verbatim}
TRAZA DE KAHN paso a paso (grafo simple):

  5 → 0 ← 4        Aristas: 5→0, 5→2, 4→0, 4→1, 2→3, 3→1
  ↓       ↓
  2 → 3 → 1

in-degree inicial:  0:2  1:2  2:1  3:1  4:0  5:0

Paso │ Queue (indeg=0) │ Saca │ Baja indeg de     │ Orden
─────┼─────────────────┼──────┼───────────────────┼──────────────
  1  │ [5, 4]          │  5   │ 0→1, 2→0 (entra)  │ 5
  2  │ [4, 2]          │  4   │ 0→0 (entra), 1→1  │ 5 4
  3  │ [2, 0]          │  2   │ 3→0 (entra)       │ 5 4 2
  4  │ [0, 3]          │  0   │ (0 no tiene salida)│ 5 4 2 0
  5  │ [3]             │  3   │ 1→0 (entra)       │ 5 4 2 0 3
  6  │ [1]             │  1   │ —                 │ 5 4 2 0 3 1

Procesados 6 = V  →  es DAG, orden válido: 5 4 2 0 3 1
\end{verbatim}
\end{dsacomofunciona}

\begin{lstlisting}
ALGORITMO DE KAHN (BFS con grados de entrada)
1. Calcula el in-degree (aristas entrantes) de cada nodo.
2. Mete a una Queue todos los nodos con in-degree = 0.
3. Mientras la Queue no esté vacía:
   - Saca un nodo u, agrégalo al resultado.
   - Por cada vecino v de u: resta 1 a su in-degree.
     Si su in-degree llega a 0, mételo a la Queue.
4. Si el resultado tiene menos nodos que V -> HAY CICLO (no es DAG).

PSEUDOCÓDIGO (Kahn)
──────────────────────────────────────────────────────────────

función topoSort(V, adj):
    indeg = arreglo de ceros de tamaño V
    PARA cada u, PARA cada v en adj[u]: indeg[v]++
    q = Queue con todos los nodos donde indeg == 0
    orden = []
    MIENTRAS q NO vacía:
        u = q.dequeue(); orden.push(u)
        PARA cada vecino v de u:
            indeg[v]--
            SI indeg[v] == 0: q.enqueue(v)
    SI orden.size < V: return "CICLO DETECTADO"
    return orden

CÓDIGO C++ (Kahn)
──────────────────────────────────────────────────────────────

#include <iostream>
#include <vector>
#include <queue>
using namespace std;

vector<int> topoSort(int V, vector<vector<int>>& adj) {
    vector<int> indeg(V, 0);
    for (int u = 0; u < V; u++)
        for (int v : adj[u]) indeg[v]++;

    queue<int> q;
    for (int i = 0; i < V; i++)
        if (indeg[i] == 0) q.push(i);

    vector<int> orden;
    while (!q.empty()) {
        int u = q.front(); q.pop();
        orden.push_back(u);
        for (int v : adj[u])
            if (--indeg[v] == 0) q.push(v);
    }
    // Si no salieron todos, había un ciclo
    if ((int)orden.size() != V) return {}; // vacío = no es DAG
    return orden;
}

CÓDIGO C++ (DFS, post-orden invertido)
──────────────────────────────────────────────────────────────

void dfs(int u, vector<vector<int>>& adj,
         vector<bool>& vis, vector<int>& pila) {
    vis[u] = true;
    for (int v : adj[u])
        if (!vis[v]) dfs(v, adj, vis, pila);
    pila.push_back(u);       // se agrega al TERMINAR
}
// El orden topológico es 'pila' invertida.

COMPLEJIDAD (Big-O)
──────────────────────────────────────────────────────────────

  Método      │ Tiempo      │ Espacio
  ────────────┼─────────────┼──────────────
  Kahn (BFS)  │ O(V + E)    │ O(V)
  DFS         │ O(V + E)    │ O(V)

* V = Vértices, E = Aristas.
\end{lstlisting}

\begin{dsaentrevistas}

✅ Detección de ciclos gratis: Kahn te dice si hay ciclo sin esfuerzo
   extra: si procesas menos de V nodos, el grafo NO es un DAG.

💡 "Course Schedule" (LeetCode 207/210): el problema estrella. "¿Puedes
   tomar todas las materias dado un mapa de prerrequisitos?" = ¿existe
   orden topológico? Respuesta directa con Kahn.

⚠️ Solo en DAGs dirigidos: si el grafo es no dirigido o tiene ciclos,
   el ordenamiento topológico NO existe. Menciónalo siempre.

🔑 Kahn vs DFS: Kahn es iterativo (sin riesgo de stack overflow) y
   detecta ciclos naturalmente. DFS es más corto pero recuerda invertir
   el post-orden al final.

\end{dsaentrevistas}

\begin{dsaresumen}

• Ordena un DAG respetando dependencias: u antes que v si u → v.
• Solo existe en grafos dirigidos SIN ciclos.
• Kahn: BFS con in-degree = 0; detecta ciclos si sobran nodos.
• DFS: apila en post-orden e invierte el resultado.
• Ambos O(V + E).
• Usos: prerrequisitos, build systems, planificación de tareas.
════════════════════════════════════════════════════════════════

\end{dsaresumen}
